\documentclass[11pt]{amsart}

\usepackage[margin=1.1in]{geometry}
\usepackage{amsmath,amssymb,amsthm,mathtools}
\usepackage{enumitem}
\usepackage{aliascnt}
\usepackage{microtype}
\usepackage{hyperref}
\usepackage[nameinlink,capitalize]{cleveref}

\hypersetup{colorlinks=true,linkcolor=blue,citecolor=blue,urlcolor=blue}
\raggedbottom

\theoremstyle{plain}
\newtheorem{theorem}{Theorem}[section]

\newaliascnt{lemma}{theorem}
\newtheorem{lemma}[lemma]{Lemma}
\aliascntresetthe{lemma}

\newaliascnt{proposition}{theorem}
\newtheorem{proposition}[proposition]{Proposition}
\aliascntresetthe{proposition}

\newaliascnt{corollary}{theorem}
\newtheorem{corollary}[corollary]{Corollary}
\aliascntresetthe{corollary}

\theoremstyle{definition}
\newaliascnt{definition}{theorem}
\newtheorem{definition}[definition]{Definition}
\aliascntresetthe{definition}


\theoremstyle{remark}
\newaliascnt{remark}{theorem}
\newtheorem{remark}[remark]{Remark}
\aliascntresetthe{remark}

\crefname{theorem}{Theorem}{Theorems}
\Crefname{theorem}{Theorem}{Theorems}
\crefname{lemma}{Lemma}{Lemmas}
\Crefname{lemma}{Lemma}{Lemmas}
\crefname{proposition}{Proposition}{Propositions}
\Crefname{proposition}{Proposition}{Propositions}
\crefname{corollary}{Corollary}{Corollaries}
\Crefname{corollary}{Corollary}{Corollaries}
\crefname{definition}{Definition}{Definitions}
\Crefname{definition}{Definition}{Definitions}
\crefname{remark}{Remark}{Remarks}
\Crefname{remark}{Remark}{Remarks}
\title[Local Concentration for Navier--Stokes]
{Local Scale-Invariant Concentration and the Type I/Type II Dichotomy
for Navier--Stokes}
\author{Guillem Duran Ballester}
\email{guillem@fragile.tech}
\subjclass[2020]{Primary 35Q30; Secondary 76D05, 35B44, 35B65}
\keywords{Navier--Stokes equations, suitable weak solutions,
\(\varepsilon\)-regularity, Type I blow-up, Type II blow-up}
\date{}

\begin{document}

\begin{abstract}
We prove a local analytic reduction for possible finite-time singularities of
the three-dimensional Navier--Stokes equations.  For a suitable weak solution,
every singular point at time \(T\) has arbitrarily small backward parabolic
cylinders on which both the standard mean-subtracted
Caffarelli--Kohn--Nirenberg scale-invariant quantity \(C+D\) and the
scale-invariant velocity quantity \(C\) are bounded below by positive
constants.  If these local scale-invariant quantities tend to zero at each
point of the singular set, the solution is regular at time \(T\) by the
\(\varepsilon\)-regularity theorem.  Along any such positive concentration
sequence, the solution either satisfies a local Type I scale-invariant bound or
falls under the local Type II alternative.
For local pointwise Type I terminal singularities of suitable weak solutions,
the endpoint weak-Serrin bridge in the companion residual paper gives the
velocity no-escape condition used in the Type I blow-up analysis.
\end{abstract}

\maketitle

\section{Introduction}

Let \((u,p)\) be a suitable weak solution of the three-dimensional
incompressible Navier--Stokes equations
\[
  \partial_t u+u\cdot\nabla u+\nabla p=\Delta u,
  \qquad \nabla\cdot u=0.
\]
We use the standard normalization in which the viscosity is one.  The weak
solution theory goes back to Leray \cite{Leray1934}.  Suitable weak solutions
and their partial regularity theory were introduced and developed by Scheffer
\cite{Scheffer1976} and by Caffarelli, Kohn, and Nirenberg \cite{CKN1982}; see
also Lin \cite{Lin1998}.  This paper establishes a local reduction from a
possible singularity at time \(T\) to the Type I/Type II blow-up rate
dichotomy.

Scale-invariant regularity and blow-up criteria provide the standard background
for this formulation; see Serrin \cite{Serrin1962}, Giga \cite{Giga1986},
Escauriaza--Seregin--{\v S}verak \cite{ESS2003}, and Seregin
\cite{Seregin2012}.  The principal regularity input used below is the
Caffarelli--Kohn--Nirenberg \(\varepsilon\)-regularity theorem.

The proof proceeds in four steps.  First, the local scale-invariant
quantities \(C\) and \(D\) are finite for suitable weak solutions and invariant
under the usual parabolic scaling.  Second, the Caffarelli--Kohn--Nirenberg
\(\varepsilon\)-regularity theorem implies that vanishing of \(C+D\) at a
candidate singular point gives local regularity, while the velocity
\(\varepsilon\)-regularity criterion gives positive concentration in the
scale-invariant \(L^3_{x,t}\) quantity.  Third, a finite-time singularity
therefore produces sequences of shrinking cylinders on which \(C+D\), and
separately \(C\), are bounded below by positive constants.  Finally, once such
a sequence is fixed, it either satisfies a Type I scale-invariant bound or
falls under the local Type II alternative.

The concentration and Type I/Type II dichotomy arguments are local in
spacetime.  Apart from the standard Caffarelli--Kohn--Nirenberg
\(\varepsilon\)-regularity theorem, all quantities and alternatives used in
that local argument are defined below.  The terminal bridge section records how
the companion residual paper upgrades the local pointwise Type I envelope to
terminal velocity no-escape through the endpoint weak-Serrin local Type I
package.

\section{Suitable Weak Solutions and Singular Points}

\begin{definition}[Suitable weak solution]
Let \(I\subset\mathbb R\) be an open interval.  A pair \((u,p)\) is a suitable
weak solution on \(\mathbb R^3\times I\) if
\[
  u\in L^\infty_tL^2_{x,\mathrm{loc}}
  \cap L^2_tH^1_{x,\mathrm{loc}},
  \qquad
  p\in L^{3/2}_{\mathrm{loc}},
  \qquad
  \nabla\cdot u=0,
\]
the Navier--Stokes equations hold in the sense of distributions, and the local
energy inequality holds: for almost every \(t_1<t_2\) in \(I\) and every
nonnegative \(\phi\in C_c^\infty(\mathbb R^3\times I)\),
\[
\begin{aligned}
&\int_{\mathbb R^3}\frac{|u(x,t_2)|^2}{2}\phi(x,t_2)\,dx
 +\int_{t_1}^{t_2}\int_{\mathbb R^3}|\nabla u|^2\phi\,dx\,dt \\
&\le
 \int_{\mathbb R^3}\frac{|u(x,t_1)|^2}{2}\phi(x,t_1)\,dx
 +\int_{t_1}^{t_2}\int_{\mathbb R^3}
   \frac{|u|^2}{2}(\partial_t\phi+\Delta\phi)\,dx\,dt \\
&\quad
 +\int_{t_1}^{t_2}\int_{\mathbb R^3}
   \left(\frac{|u|^2}{2}+p\right)u\cdot\nabla\phi\,dx\,dt .
\end{aligned}
\]
The pressure is understood modulo functions of time.
\end{definition}

For \(z_0=(x_0,t_0)\) and \(r>0\), write
\[
  Q_r(z_0)=B_r(x_0)\times(t_0-r^2,t_0).
\]
We also write \(Q_r(x_0,t_0)\) for \(Q_r((x_0,t_0))\).

\begin{definition}[Singular set at a fixed time]
For a time \(T\), define
\[
  \Sigma(T)=
  \{x\in\mathbb R^3:\ u\text{ is not locally bounded in any }
  Q_r(x,T)\}.
\]
Equivalently, \(x\notin\Sigma(T)\) if there are \(r>0\) and \(M<\infty\) such
that \(\|u\|_{L^\infty(Q_r(x,T))}\le M\).
\end{definition}

If a finite-time singularity occurs at time \(T\), then
\(\Sigma(T)\ne\emptyset\) in this local sense.

\section{Parabolic Rescaling and Scale-Invariant Quantities}

Fix \(z_0=(x_0,T)\).  For \(r>0\), define the parabolic rescaling
\[
  u^{(r)}(y,s)=r\,u(x_0+ry,T+r^2s),
  \qquad
  p^{(r)}(y,s)=r^2p(x_0+ry,T+r^2s).
\]
The rescaled variables are defined on the set of \((y,s)\) for which
\((x_0+ry,T+r^2s)\) belongs to the original spacetime domain.

\begin{proposition}[Invariance under parabolic rescaling]
If \((u,p)\) is a suitable weak solution near \((x_0,T)\), then
\((u^{(r)},p^{(r)})\) is a suitable weak solution on the rescaled domain.  The
local energy inequality is preserved, and the pressure remains defined modulo
functions of the rescaled time.
\end{proposition}

\begin{proof}
The distributional equations and the divergence condition follow from the
change of variables \(x=x_0+ry\), \(t=T+r^2s\).  We include the verification of
the local energy inequality to fix the normalization.

Let \(\psi\in C_c^\infty\) be nonnegative and compactly supported in the
rescaled domain.  Define
\[
  \phi(x,t)=
  \psi\left(\frac{x-x_0}{r},\frac{t-T}{r^2}\right).
\]
If \(y=(x-x_0)/r\) and \(s=(t-T)/r^2\), then
\[
  \partial_t\phi=r^{-2}(\partial_s\psi)(y,s),\qquad
  \nabla_x\phi=r^{-1}(\nabla_y\psi)(y,s),\qquad
  \Delta_x\phi=r^{-2}(\Delta_y\psi)(y,s).
\]
For almost every \(s_1<s_2\), the corresponding times
\(T+r^2s_1<T+r^2s_2\) are admissible in the original local energy inequality.
Apply that inequality for \((u,p)\) to \(\phi\) on
\([T+r^2s_1,T+r^2s_2]\).  After the change of variables
\(x=x_0+ry\), \(t=T+r^2s\), using
\[
  u(x_0+ry,T+r^2s)=r^{-1}u^{(r)}(y,s),\qquad
  p(x_0+ry,T+r^2s)=r^{-2}p^{(r)}(y,s),
\]
each term contains the same common factor \(r\).  Dividing by this factor gives
the rescaled inequality
\[
\begin{aligned}
&\int_{\mathbb R^3}\frac{|u^{(r)}(y,s_2)|^2}{2}\psi(y,s_2)\,dy
 +\int_{s_1}^{s_2}\int_{\mathbb R^3}|\nabla_y u^{(r)}|^2\psi\,dy\,ds \\
&\le
 \int_{\mathbb R^3}\frac{|u^{(r)}(y,s_1)|^2}{2}\psi(y,s_1)\,dy
 +\int_{s_1}^{s_2}\int_{\mathbb R^3}
   \frac{|u^{(r)}|^2}{2}(\partial_s\psi+\Delta_y\psi)\,dy\,ds \\
&\quad
 +\int_{s_1}^{s_2}\int_{\mathbb R^3}
   \left(\frac{|u^{(r)}|^2}{2}+p^{(r)}\right)
   u^{(r)}\cdot\nabla_y\psi\,dy\,ds .
\end{aligned}
\]
Thus \((u^{(r)},p^{(r)})\) is suitable on the rescaled domain.  The pressure
transforms according to the Navier--Stokes scaling, and addition of a function
of \(t\) becomes addition of a function of \(s\).
\end{proof}

\begin{definition}[Scale-invariant local quantities]
For \(z_0=(x_0,t_0)\) and \(r>0\), set
\[
  C(z_0,r)=r^{-2}\int_{Q_r(z_0)} |u|^3\,dx\,dt
\]
and
\[
  D(z_0,r)=r^{-2}
  \int_{t_0-r^2}^{t_0}\int_{B_r(x_0)}
  |p(x,t)-(p)_{B_r(x_0)}(t)|^{3/2}\,dx\,dt.
\]
Here \((p)_{B_r(x_0)}(t)\) denotes the spatial mean of \(p(\cdot,t)\) on
\(B_r(x_0)\).
\end{definition}

The subtraction of the spatial mean removes the pressure ambiguity in the form
used by the Caffarelli--Kohn--Nirenberg criterion.

\begin{proposition}[Local finiteness and scaling]\label{prop:critical-scaling}
Let \((u,p)\) be a suitable weak solution near \(z_0=(x_0,T)\).  Then
\(C(z_0,r)\) and \(D(z_0,r)\) are finite for every sufficiently small \(r\).
Moreover, for every fixed \(R<\infty\),
\[
\begin{aligned}
&\int_{-R^2}^{0}\int_{B_R}
\left(
|u^{(r)}|^3+
|p^{(r)}-(p^{(r)})_{B_R}(s)|^{3/2}
\right)\,dy\,ds \\
&\qquad
=R^2\bigl(C(z_0,rR)+D(z_0,rR)\bigr).
\end{aligned}
\]
\end{proposition}

\begin{proof}
Local finiteness of the velocity term follows from
\[
  u\in L^\infty_tL^2_{x,\mathrm{loc}}
  \cap L^2_tH^1_{x,\mathrm{loc}}
  \subset L^3_{\mathrm{loc}}
\]
on finite cylinders.  Indeed, if \(B\Subset\mathbb R^3\) and
\((a,b)\Subset I\), Sobolev and interpolation give
\[
  \int_a^b \|u(t)\|_{L^3(B)}^3\,dt
  \le
  C
  \left(\operatorname*{ess\,sup}_{a<t<b}\|u(t)\|_{L^2(B)}\right)^{3/2}
  \left(\int_a^b \|u(t)\|_{H^1(B)}^2\,dt\right)^{3/4}
  (b-a)^{1/4}.
\]
Thus \(u\in L^3_{\mathrm{loc}}\).  The pressure term is finite because
\(p\in L^{3/2}_{\mathrm{loc}}\), and subtracting a spatial mean over the same
ball preserves local \(L^{3/2}\) integrability.  More explicitly, for a ball
\(B\),
\[
  |p-(p)_B(t)|^{3/2}\le C\left(|p|^{3/2}+|(p)_B(t)|^{3/2}\right),
\]
while Jensen's inequality gives
\[
  |(p)_B(t)|^{3/2}
  \le \frac1{|B|}\int_B |p(x,t)|^{3/2}\,dx.
\]
After integration over \(B\) and over time, the mean-subtracted pressure term is
controlled by the local \(L^{3/2}\) norm of \(p\).

For the scaling identity, use \(x=x_0+ry\), \(t=T+r^2s\).  The velocity term
transforms as
\[
  \int_{-R^2}^{0}\int_{B_R}|u^{(r)}|^3\,dy\,ds
  =r^{-2}\int_{T-(rR)^2}^{T}\int_{B_{rR}(x_0)}|u|^3\,dx\,dt.
\]
The pressure means transform according to
\[
  (p^{(r)})_{B_R}(s)=r^2(p)_{B_{rR}(x_0)}(T+r^2s),
\]
because
\[
  \frac1{|B_R|}\int_{B_R} r^2p(x_0+ry,T+r^2s)\,dy
  =
  \frac{r^2}{|B_{rR}|}\int_{B_{rR}(x_0)}p(x,T+r^2s)\,dx .
\]
Therefore
\[
\begin{aligned}
&\int_{-R^2}^{0}\int_{B_R}
|p^{(r)}-(p^{(r)})_{B_R}(s)|^{3/2}\,dy\,ds \\
&\qquad =
r^{-2}\int_{T-(rR)^2}^{T}\int_{B_{rR}(x_0)}
|p-(p)_{B_{rR}(x_0)}(t)|^{3/2}\,dx\,dt .
\end{aligned}
\]
Adding the velocity and pressure identities gives the displayed formula.
\end{proof}

\section{\texorpdfstring{\(\varepsilon\)-Regularity}{Epsilon-Regularity} and the Vanishing Case}

\begin{theorem}[Caffarelli--Kohn--Nirenberg \(\varepsilon\)-regularity]
\label{thm:ckn}
There exists a universal constant \(\varepsilon_0>0\) with the
following property.  Let \((u,p)\) be a suitable weak solution in \(Q_r(z_0)\).
If
\[
  C(z_0,r)+D(z_0,r)<\varepsilon_0,
\]
then \(u\) is locally bounded in a smaller cylinder, for instance in
\(Q_{r/2}(z_0)\).  In particular, if \(z_0=(x_0,T)\), then
\(x_0\notin\Sigma(T)\).
\end{theorem}

\begin{proof}
This is the local \(\varepsilon\)-regularity theorem of
Caffarelli--Kohn--Nirenberg \cite{CKN1982}; see also Lin \cite{Lin1998}.  It is
stated here with the spatial mean subtraction used in \(D\).  Equivalent
representatives of the pressure differ by functions of time, and the spatial
mean subtraction on \(B_r(x_0)\) removes this ambiguity.
\end{proof}

\begin{theorem}[Velocity \(\varepsilon\)-regularity]\label{thm:velocity-ckn}
There exists a universal constant \(\varepsilon_v>0\) with the following
property.  Let \((u,p)\) be a suitable weak solution in \(Q_r(z_0)\).  If
\[
  C(z_0,r)<\varepsilon_v,
\]
then \(u\) is locally bounded in \(Q_{r/2}(z_0)\).  In particular, if
\(z_0=(x_0,T)\), then \(x_0\notin\Sigma(T)\).
\end{theorem}

\begin{proof}
This is the standard \(L^3_{x,t}\) form of the Caffarelli--Kohn--Nirenberg
\(\varepsilon\)-regularity criterion.  It follows from the standard local
pressure decay lemma together with \cref{thm:ckn}; equivalently, it is the
\(L^3_{x,t}\)-smallness criterion
proved in the Caffarelli--Kohn--Nirenberg theory and in Lin's direct proof
\cite{CKN1982,Lin1998}.  The statement is scale invariant because \(C\) is
invariant under the parabolic Navier--Stokes scaling.
\end{proof}

\begin{definition}[Vanishing local scale-invariant quantities at time \(T\)]
We say that the local scale-invariant quantities vanish at time \(T\) if, for
every \(x_0\in\Sigma(T)\),
\[
  \limsup_{r\downarrow0}
  \bigl(C((x_0,T),r)+D((x_0,T),r)\bigr)=0.
\]
This condition is vacuous when \(\Sigma(T)=\emptyset\).
\end{definition}

\begin{proposition}[Equivalent rescaled formulation]
\label{prop:vanishing-rescaled}
The preceding condition is equivalent to the following statement: for every
\(x_0\in\Sigma(T)\) and every fixed \(R<\infty\),
\[
  \lim_{r\downarrow0}
  \int_{-R^2}^{0}\int_{B_R}
  \left(
  |u^{(r)}|^3+
  |p^{(r)}-(p^{(r)})_{B_R}(s)|^{3/2}
  \right)\,dy\,ds=0.
\]
\end{proposition}

\begin{proof}
By \cref{prop:critical-scaling}, the rescaled integral over
\(B_R\times(-R^2,0)\) equals
\[
  R^2\bigl(C((x_0,T),rR)+D((x_0,T),rR)\bigr).
\]
Thus vanishing of \(C+D\) implies vanishing on every fixed rescaled cylinder.
Conversely, taking \(R=1\) gives the original pointwise condition.
\end{proof}

\begin{theorem}[Regularity under vanishing local quantities]\label{thm:no-concentration}
Let \((u,p)\) be a suitable weak solution on
\(\mathbb R^3\times(T-\delta,T)\) for some \(\delta>0\).  If the local
scale-invariant quantities vanish at time \(T\), then
\(\Sigma(T)=\emptyset\).  Thus no finite-time singularity occurs at time \(T\)
in the no-concentration case.
\end{theorem}

\begin{proof}
Suppose, toward a contradiction, that \(x_0\in\Sigma(T)\).  By the vanishing
hypothesis, choose \(r>0\) such that
\[
  C((x_0,T),r)+D((x_0,T),r)<\varepsilon_0.
\]
\Cref{thm:ckn} implies that \(u\) is bounded in \(Q_{r/2}(x_0,T)\), hence
\(x_0\notin\Sigma(T)\).  This contradicts the choice of \(x_0\).  Hence
\(\Sigma(T)=\emptyset\).
\end{proof}

\begin{corollary}[Vanishing local quantities exclude a terminal singularity]
\label{cor:escape}
Let \((u,p)\) be suitable near \((x_0,T)\).  If
\[
  \limsup_{r\downarrow0}
  \bigl(C((x_0,T),r)+D((x_0,T),r)\bigr)=0,
\]
then \(u\) is bounded in some backward cylinder \(Q_\rho(x_0,T)\), and hence
\(x_0\notin\Sigma(T)\).  Equivalently, membership in \(\Sigma(T)\) requires
that \(C+D\) have a positive lower bound along a sequence of shrinking backward
parabolic cylinders.
\end{corollary}

\begin{proof}
The hypothesis gives a scale at which \(C+D\) is below the universal threshold
in \cref{thm:ckn}.  The \(\varepsilon\)-regularity theorem gives local boundedness in
a backward cylinder ending at \(T\), so \(x_0\notin\Sigma(T)\).
\end{proof}

\section{Positive Local Concentration}

\begin{lemma}[Failure of vanishing gives a positive concentration sequence]
\label{lem:positive-concentration}
Suppose the local scale-invariant quantities do not vanish at time \(T\).  Then
there exist \(x_0\in\Sigma(T)\), a constant \(\eta>0\), and radii
\(r_n\downarrow0\) such that
\[
  C((x_0,T),r_n)+D((x_0,T),r_n)\ge\eta
  \qquad\text{for all }n.
\]
Equivalently, the rescaled sequence around \((x_0,T)\) has a positive lower
bound for the corresponding local scale-invariant quantities on a fixed
bounded cylinder.
\end{lemma}

\begin{proof}
Negating the definition of vanishing gives a point \(x_0\in\Sigma(T)\) with
positive limsup of
\[
  C((x_0,T),r)+D((x_0,T),r)
\]
as \(r\downarrow0\).  Choose \(\eta>0\) below that limsup and choose a sequence
\(r_n\downarrow0\) along which the displayed lower bound holds.  The rescaled
formulation follows from \cref{prop:critical-scaling}.
\end{proof}

\begin{definition}[Positive local concentration sequence]
\label{def:positive-concentration-sequence}
Let \(x_0\in\Sigma(T)\).  A sequence \(r_n\downarrow0\) is a positive local
concentration sequence at \((x_0,T)\) if there is \(\eta>0\) such that
\[
  C((x_0,T),r_n)+D((x_0,T),r_n)\ge\eta
  \qquad\text{for all }n.
\]
The corresponding parabolic rescalings are those defined by
\[
  u_n(y,s)=r_nu(x_0+r_ny,T+r_n^2s),
  \qquad
  p_n(y,s)=r_n^2p(x_0+r_ny,T+r_n^2s).
\]
\end{definition}

\begin{proposition}[Persistence of positive concentration]
\label{prop:compactness-persistence}
Let \((u_n,p_n)\) be suitable weak solutions on a fixed cylinder
\(B_R\times(-R^2,0)\).  Suppose that
\[
  \int_{-R^2}^{0}\int_{B_R}
  \left(|u_n|^3+|p_n-(p_n)_{B_R}(s)|^{3/2}\right)\,dy\,ds
  \ge\eta>0
\]
for all \(n\).  Suppose further that, after passing to a subsequence,
\[
  u_n\to U\quad\text{in }L^3(B_R\times(-R^2,0))
\]
and
\[
  p_n-(p_n)_{B_R}(s)\to \Pi
  \quad\text{in }L^{3/2}(B_R\times(-R^2,0)).
\]
Then
\[
  \int_{-R^2}^{0}\int_{B_R}
  \left(|U|^3+|\Pi|^{3/2}\right)\,dy\,ds\ge\eta.
\]
In particular, the limiting pair retains the same positive lower bound for the
mean-subtracted local scale-invariant quantity.  If
\(\Pi=P-(P)_{B_R}(s)\), this is the corresponding
Caffarelli--Kohn--Nirenberg quantity of \((U,P)\) on
\(B_R\times(-R^2,0)\).
\end{proposition}

\begin{proof}
Strong convergence in \(L^3\) gives
\[
  \int_{-R^2}^{0}\int_{B_R}|u_n|^3\,dy\,ds
  \longrightarrow
  \int_{-R^2}^{0}\int_{B_R}|U|^3\,dy\,ds.
\]
Strong convergence in \(L^{3/2}\) gives the analogous convergence of the
mean-subtracted pressure integrals.  Passing to the limit in the lower bound
yields the asserted inequality.
\end{proof}

\begin{remark}
\Cref{prop:compactness-persistence} is an auxiliary compactness observation,
independent of the proof of the dichotomy below.
\end{remark}

\begin{proposition}[Local alternative]\label{prop:local-alternative}
Let \((u,p)\) be a suitable weak solution on
\(\mathbb R^3\times(T-\delta,T)\) for some \(\delta>0\).  Exactly one of the
following alternatives holds.
\begin{enumerate}[label=\textup{(\roman*)}]
\item For every \(x_0\in\Sigma(T)\),
\[
  \limsup_{r\downarrow0}
  \bigl(C((x_0,T),r)+D((x_0,T),r)\bigr)=0.
\]
\item There exist \(x_0\in\Sigma(T)\), \(\eta>0\), and \(r_n\downarrow0\) such
that
\[
  C((x_0,T),r_n)+D((x_0,T),r_n)\ge\eta
  \qquad\text{for all }n.
\]
\end{enumerate}
\end{proposition}

\begin{proof}
The first alternative is a universal assertion over the singular set, and the
second is its failure expressed by a sequence.  If the first alternative fails,
then for some \(x_0\in\Sigma(T)\) the corresponding limsup is positive;
choosing a positive number below this limsup and a sequence along which the
lower bound holds gives the second alternative.  Conversely, the second
alternative contradicts the first.  Hence exactly one alternative holds.
\end{proof}

\begin{theorem}[Finite-time singularity yields positive local scale-invariant concentration]
\label{thm:singular-positive}
Let \((u,p)\) be a suitable weak solution on
\(\mathbb R^3\times(T-\delta,T)\).  If \(\Sigma(T)\ne\emptyset\), then there
are \(x_0\in\Sigma(T)\), \(\eta>0\), and \(r_n\downarrow0\) such that
\[
  C((x_0,T),r_n)+D((x_0,T),r_n)\ge\eta
  \qquad\text{for all }n.
\]
\end{theorem}

\begin{proof}
If the first alternative in \cref{prop:local-alternative} held, then
\cref{thm:no-concentration} would imply \(\Sigma(T)=\emptyset\), contrary to
the hypothesis.  Therefore the second alternative in
\cref{prop:local-alternative} holds.
\end{proof}

\begin{corollary}[Singular points carry positive velocity concentration]
\label{cor:velocity-concentration}
Let \((u,p)\) be a suitable weak solution on
\(\mathbb R^3\times(T-\delta,T)\).  If \(x_0\in\Sigma(T)\), then there are
\(\eta_v>0\) and \(r_n\downarrow0\) such that
\[
  C((x_0,T),r_n)\ge\eta_v
  \qquad\text{for all }n.
\]
\end{corollary}

\begin{proof}
If the conclusion failed, then
\[
  \limsup_{r\downarrow0} C((x_0,T),r)=0.
\]
Choose \(r>0\) so small that
\[
  C((x_0,T),r)<\varepsilon_v,
\]
where \(\varepsilon_v\) is the constant in
\cref{thm:velocity-ckn}.  The velocity \(\varepsilon\)-regularity theorem gives
boundedness of \(u\) in \(Q_{r/2}(x_0,T)\), contradicting
\(x_0\in\Sigma(T)\).  Hence the limsup is positive; taking any
\(\eta_v\) below it and passing to a sequence with
\(C((x_0,T),r_n)\ge\eta_v\) gives the result.
\end{proof}

\section{Automatic No-Escape for Finite-Energy Type I Singularities}
\label{sec:auto-no-escape}

The preceding concentration result gives the positive velocity concentration
sequence used in the Type I blow-up analysis.  The remaining compactness input
is a no-escape condition excluding the possibility that all of the selected
scale-invariant velocity mass collapses into the final rescaled time slice.
We prove here that this no-escape condition is automatic for finite-energy
Type I singularities.

For this section we use, in addition to \(C\) and \(D\), the standard local
kinetic-energy and dissipation quantities
\[
  A(z_0,r)
  =r^{-1}\operatorname*{ess\,sup}_{t_0-r^2<t<t_0}
    \int_{B_r(x_0)} |u(x,t)|^2\,dx,
\]
and
\[
  E(z_0,r)
  =r^{-1}\int_{t_0-r^2}^{t_0}\int_{B_r(x_0)}
    |\nabla u(x,t)|^2\,dx\,dt .
\]
These quantities are invariant under the Navier--Stokes scaling.

\begin{definition}[Finite-energy Type I terminal singularity]
\label{def:finite-energy-typeI-paper0}
Let \((u,p)\) be a suitable weak solution on \(\mathbb R^3\times(0,T)\).
We say that \(z_*=(x_*,T)\) is a finite-energy Type I terminal singularity if
\(x_*\in\Sigma(T)\),
\[
  u\in L^\infty(0,T;L^2(\mathbb R^3))
  \cap L^2(0,T;\dot H^1(\mathbb R^3)),
\]
and there are \(T_0<T\) and \(M<\infty\) such that
\begin{equation}\label{eq:paper0-global-typeI}
  \|u(t)\|_{L^\infty(\mathbb R^3)}
  \le \frac{M}{\sqrt{T-t}},
  \qquad T_0<t<T .
\end{equation}
\end{definition}

The key point is a uniformly-local energy propagation estimate.  It is stated
at unit scale, because the application is obtained by parabolic rescaling from
\((x_*,T)\).  The estimate is a local state-space closure: the state is the
uniformly-local kinetic energy, and the Type I envelope turns the nonlinear
pressure stratum into an integrable linear coefficient.  No bound below
depends on the size of the global \(L^2\) energy or the global dissipation;
finite energy is used only to place the pressure in the standard whole-space
Leray gauge and to make fixed-scale terminal \(L^3\) absolute continuity
immediate.

We use the standard Leray pressure gauge throughout this section.  At almost
every terminal time on which the Type I bound holds,
\(u(t)\in L^2(\mathbb R^3)\cap L^\infty(\mathbb R^3)\), hence
\(u_i(t)u_j(t)\in L^{3/2}(\mathbb R^3)\).  The whole-space pressure
\[
  p^L(t)=\mathcal R_i\mathcal R_j(u_i(t)u_j(t))
\]
is then defined by the Calderon--Zygmund theory.  The distributional pressure
equation gives
\(-\Delta p=\partial_i\partial_j(u_i u_j)=-\Delta p^L\), and the whole-space
de Rham argument applied to the momentum equation shows that \(p-p^L\) is a
function of time.  Thus replacing \(p\) by \(p^L\) is only a pressure-gauge
change and does not alter suitability.  This gauge choice uses the
finite-energy setting to define the representative; the uniform constants
below, in particular \(B_{\mathrm{uloc}}\) and \(B_A\), do not contain
\(\|u\|_{L^\infty(0,T;L^2(\mathbb R^3))}\) or
\(\int_0^T\|\nabla u(t)\|_{L^2(\mathbb R^3)}^2\,dt\).

\begin{lemma}[Uniformly-local energy propagation under a Type I envelope]
\label{lem:uloc-typeI-propagation}
Let \((v,\pi)\) be a suitable weak solution on
\(\mathbb R^3\times(-4,0)\) with \(v(s)\in L^2(\mathbb R^3)\) for a.e.
\(s\), and take \(\pi\) to be the whole-space Leray pressure representative,
modulo functions of time.  Assume that
\begin{equation}\label{eq:unit-typeI-envelope}
  |v(x,s)|\le M(-s)^{-1/2},
  \qquad \text{for a.e. }(x,s)\in\mathbb R^3\times(-4,0),
\end{equation}
and that for some admissible initial time \(s_0\in(-4,-3)\),
\begin{equation}\label{eq:unit-uloc-initial}
  \alpha_0
  :=\sup_{a\in\mathbb R^3}\int_{B_2(a)} |v(x,s_0)|^2\,dx
  \le K .
\end{equation}
Then there is a constant \(B_{\mathrm{uloc}}=B_{\mathrm{uloc}}(M,K)<\infty\)
such that
\begin{equation}\label{eq:unit-uloc-bound}
  \operatorname*{ess\,sup}_{-1<s<0}\sup_{a\in\mathbb R^3}
  \int_{B_1(a)} |v(x,s)|^2\,dx
  +
  \sup_{a\in\mathbb R^3}\int_{-1}^{0}\int_{B_1(a)}
  |\nabla v|^2\,dx\,ds
  \le B_{\mathrm{uloc}} .
\end{equation}
\end{lemma}

\begin{proof}
We give the proof because the estimate is the mechanism that closes the
terminal no-escape gap.  Let
\[
  \alpha(s)=\sup_{a\in\mathbb R^3}\int_{B_1(a)} |v(x,s)|^2\,dx .
\]
By \eqref{eq:unit-typeI-envelope},
\begin{equation}\label{eq:alpha-typeI-trivial}
  \alpha(s)^{1/2}\le C M(-s)^{-1/2},
  \qquad -4<s<0 .
\end{equation}
Fix \(a\in\mathbb R^3\) and choose
\(\phi_a\in C_c^\infty(B_2(a))\) with
\(0\le\phi_a\le1\), \(\phi_a\equiv1\) on \(B_1(a)\), and
\(|\nabla\phi_a|+|\Delta\phi_a|\le C\), uniformly in \(a\).  Applying the
local energy inequality on \((s_0,t)\), using the standard monotone
approximation of temporal cutoffs and with the pressure gauge chosen below,
gives for a.e. \(t\in(s_0,0)\)
\begin{align}
 &\int |v(t)|^2\phi_a\,dx
   +2\int_{s_0}^{t}\int |\nabla v|^2\phi_a\,dx\,ds       \notag \\
 &\quad\le
   \int |v(s_0)|^2\phi_a\,dx
   +C\int_{s_0}^{t}\int_{B_2(a)} |v|^2\,dx\,ds
   +C\int_{s_0}^{t}\int_{B_2(a)} |v|^3\,dx\,ds          \notag \\
 &\qquad
   +C\int_{s_0}^{t}\int_{B_2(a)}
       |\pi-c_a(s)|\,|v|\,dx\,ds .
\label{eq:lei-unit-estimate}
\end{align}
Here \(c_a(s)\) is an arbitrary function of time; adding or subtracting such a
function does not change the local energy inequality, since
\(\nabla\cdot v=0\) and \(\phi_a\) is compactly supported.

The velocity terms are controlled by the state.  Since \(B_2(a)\) is covered by
finitely many unit balls,
\[
  \int_{B_2(a)} |v|^2\,dx\le C\alpha(s),
\]
and by \eqref{eq:unit-typeI-envelope},
\begin{equation}\label{eq:cubic-unit-bound}
  \int_{B_2(a)} |v|^3\,dx
  \le M(-s)^{-1/2}\int_{B_2(a)} |v|^2\,dx
  \le C M(-s)^{-1/2}\alpha(s).
\end{equation}
It remains to estimate the pressure term uniformly in \(a\).

Use the whole-space pressure representative
\(\pi=\mathcal R_i\mathcal R_j(v_i v_j)\), modulo functions of time.  Let
\(\chi_a\in C_c^\infty(B_8(a))\), \(\chi_a\equiv1\) on \(B_4(a)\), and write
inside \(B_4(a)\)
\[
  \pi=q_a+h_a,
  \qquad
  q_a=\mathcal R_i\mathcal R_j(\chi_a v_i v_j).
\]
Then \(h_a\) is harmonic in \(B_4(a)\).  Calderon--Zygmund boundedness
\cite{Stein1970} and \eqref{eq:unit-typeI-envelope} give, for a.e. \(s\),
\begin{align*}
  \|q_a(s)\|_{L^2(B_4(a))}
  &\le C\|\chi_a v(s)\otimes v(s)\|_{L^2(\mathbb R^3)}  \\
  &\le C M(-s)^{-1/2}\|v(s)\|_{L^2(B_8(a))}             \\
  &\le C M(-s)^{-1/2}\alpha(s)^{1/2},
\end{align*}
where the last step again uses a finite covering by unit balls.  Hence
\begin{equation}\label{eq:local-pressure-unit}
  \int_{B_2(a)} |q_a|\,|v|\,dx
  \le C M(-s)^{-1/2}\alpha(s).
\end{equation}

For the harmonic part, set \(c_a(s)=h_a(a,s)\).  The kernel estimate
\[
  |K_{ij}(x-z)-K_{ij}(a-z)|
  \le C |x-a|\,|z-a|^{-4},
  \qquad x\in B_2(a),\quad |z-a|>4,
\]
where \(K_{ij}\) is the kernel of \(\mathcal R_i\mathcal R_j\), gives
\[
  |h_a(x,s)-c_a(s)|
  \le
  C\sum_{m=2}^{\infty}2^{-4m}
    \int_{2^m<|z-a|<2^{m+1}} |v(z,s)|^2\,dz .
\]
The annulus \(\{2^m<|z-a|<2^{m+1}\}\) is covered by at most \(C2^{3m}\)
unit balls.  Therefore
\begin{equation}\label{eq:harmonic-pressure-unit}
  \|h_a(s)-c_a(s)\|_{L^\infty(B_2(a))}
  \le C\sum_{m=2}^{\infty}2^{-m}\alpha(s)
  \le C\alpha(s).
\end{equation}
Consequently,
\begin{align}
  \int_{B_2(a)} |h_a-c_a|\,|v|\,dx
  &\le C\alpha(s)\int_{B_2(a)} |v|\,dx                    \notag\\
  &\le C\alpha(s)^{3/2}                                    \notag\\
  &\le C M(-s)^{-1/2}\alpha(s),
\label{eq:harmonic-pressure-flux-unit}
\end{align}
where the last inequality is exactly the Type I state-space closure
\eqref{eq:alpha-typeI-trivial}.
Combining \eqref{eq:local-pressure-unit} and
\eqref{eq:harmonic-pressure-flux-unit},
\begin{equation}\label{eq:pressure-flux-unit}
  \int_{B_2(a)} |\pi-c_a(s)|\,|v|\,dx
  \le C M(-s)^{-1/2}\alpha(s).
\end{equation}

Insert \eqref{eq:cubic-unit-bound} and \eqref{eq:pressure-flux-unit} into
\eqref{eq:lei-unit-estimate}, then take the supremum over \(a\).  Since
\((-s)^{-1/2}\in L^1(-4,0)\), Gronwall's inequality yields
\[
  \operatorname*{ess\,sup}_{s_0<s<0}\alpha(s)
  \le C K
     \exp\left(C\int_{s_0}^{0}\bigl(1+M(-s)^{-1/2}\bigr)\,ds\right)
  \le B_1(M,K).
\]
Returning to \eqref{eq:lei-unit-estimate}, using this bound, and then letting
\(t\uparrow0\) through admissible times gives the corresponding
uniformly-local dissipation estimate on \((-1,0)\).  This proves
\eqref{eq:unit-uloc-bound}.
\end{proof}

\begin{remark}[State-space stratification]
\label{rem:state-space-stratification}
The preceding proof can be read as a three-stratum closure of the local state
\(\alpha(s)\).  The diffusion and transport strata are linear in
\(\alpha\), with coefficient \(1+M(-s)^{-1/2}\).  The local pressure stratum is
also linear after Calderon--Zygmund and the Type I envelope.  The only
apparently superlinear stratum is the harmonic pressure flux, which is
\(O(\alpha^{3/2})\); the pointwise Type I envelope implies
\(\alpha^{1/2}\le C M(-s)^{-1/2}\), reducing it to the same integrable linear
coefficient.  Thus the bad region \(\alpha\gg1\) is not invariant: the
state-space dynamics are trapped in a bounded absorbing interval depending
only on \(M\) and the initial uniformly-local energy state.
\end{remark}

\begin{proposition}[Automatic terminal kinetic tightness]
\label{prop:auto-terminal-A}
Let \((u,p,z_*)\), \(z_*=(x_*,T)\), be a finite-energy Type I terminal
singularity in the sense of \cref{def:finite-energy-typeI-paper0}.  Then there
exist \(r_0>0\) and \(B_A<\infty\), with \(B_A\) depending only on the Type I
constant \(M\), such that
\begin{equation}\label{eq:auto-A-bound}
  A(z_*,r)\le B_A,
  \qquad 0<r<r_0 .
\end{equation}
Moreover,
\begin{equation}\label{eq:auto-E-bound}
  r^{-1}\int_{T-r^2}^{T}\int_{B_r(x_*)}|\nabla u|^2\,dx\,dt
  \le B_A,
  \qquad 0<r<r_0 .
\end{equation}
\end{proposition}

\begin{proof}
Choose \(r_0>0\) so small that \(T-4r^2>T_0\) for \(0<r<r_0\).  For such
\(r\), use the whole-space Leray pressure gauge fixed above and define the
rescaled solution
\[
  v^{(r)}(y,s)=r u(x_*+ry,T+r^2s),
  \qquad
  \pi^{(r)}(y,s)=r^2 p(x_*+ry,T+r^2s),
\]
on \(\mathbb R^3\times(-4,0)\).  The Type I bound
\eqref{eq:paper0-global-typeI} gives
\[
  |v^{(r)}(y,s)|\le M(-s)^{-1/2},
  \qquad -4<s<0 .
\]
For any admissible \(s_0\in(-4,-3)\) at which the Type I bound holds, and any
\(a\in\mathbb R^3\),
\begin{align*}
  \int_{B_2(a)} |v^{(r)}(y,s_0)|^2\,dy
  &= r^{-1}\int_{B_{2r}(x_*+ra)}
       |u(x,T+r^2s_0)|^2\,dx                                    \\
  &\le r^{-1}|B_{2r}|\frac{M^2}{r^2|s_0|}
   \le C M^2 .
\end{align*}
Such times form a full-measure subset of \((-4,-3)\).  The preceding estimate
is uniform in this choice, so \cref{lem:uloc-typeI-propagation} applies with
\(K=CM^2\), uniformly in \(r\).  Hence
\[
  \operatorname*{ess\,sup}_{-1<s<0}\int_{B_1}|v^{(r)}(y,s)|^2\,dy
  +\int_{-1}^{0}\int_{B_1}|\nabla v^{(r)}|^2\,dy\,ds
  \le B_A .
\]
Undoing the parabolic change of variables gives exactly
\eqref{eq:auto-A-bound} and \eqref{eq:auto-E-bound}.  The constant \(B_A\)
depends only on \(M\); no global energy norm of \(u\) has been used in the
estimate.
\end{proof}

\begin{theorem}[Automatic velocity no-escape]
\label{thm:auto-velocity-no-escape}
Let \((u,p,z_*)\), \(z_*=(x_*,T)\), be a finite-energy Type I terminal
singularity.  Let \(\lambda_k\downarrow0\) be any sequence of scales, and set
\[
  u_k(x,t)=\lambda_k u(x_*+\lambda_k x,T+\lambda_k^2 t).
\]
Then
\begin{equation}\label{eq:auto-no-escape}
  \lim_{\sigma\downarrow0}
  \sup_k\int_{-\sigma}^{0}\int_{B_1}|u_k(x,t)|^3\,dx\,dt=0 .
\end{equation}
\end{theorem}

\begin{proof}
By \cref{prop:auto-terminal-A}, after discarding at most finitely many indices,
\[
  \operatorname*{ess\,sup}_{-1<t<0}\int_{B_1}|u_k(x,t)|^2\,dx\le B_A .
\]
The Type I bound gives, for the same tail of the sequence,
\[
  \|u_k(t)\|_{L^\infty(\mathbb R^3)}
  \le M(-t)^{-1/2},
  \qquad -1<t<0 .
\]
Therefore, for \(0<\sigma<1\),
\[
\begin{aligned}
  \int_{-\sigma}^{0}\int_{B_1}|u_k|^3\,dx\,dt
  &\le
  \int_{-\sigma}^{0}\|u_k(t)\|_{L^\infty(B_1)}
       \int_{B_1}|u_k(x,t)|^2\,dx\,dt  \\
  &\le M B_A\int_{-\sigma}^{0}(-t)^{-1/2}\,dt
   =2M B_A\sigma^{1/2} .
\end{aligned}
\]
This bound is uniform over the tail.  For each of the finitely many discarded
indices \(k\), choose \(\sigma_k>0\) so that
\(T+\lambda_k^2t>T_0\) for \(-\sigma_k<t<0\).  On this terminal subinterval
the Type I envelope and the finite-energy assumption imply
\[
  \int_{-\sigma_k}^{0}\int_{B_1}|u_k|^3\,dx\,dt
  \le
  \int_{-\sigma_k}^{0}\|u_k(t)\|_{L^\infty(B_1)}
       \int_{B_1}|u_k(x,t)|^2\,dx\,dt
  <\infty .
\]
Indeed, on \((-\sigma_k,0)\),
\[
  \|u_k(t)\|_{L^\infty(B_1)}\le M(-t)^{-1/2},
  \qquad
  \int_{B_1}|u_k(x,t)|^2\,dx
  \le \lambda_k^{-1}\|u(T+\lambda_k^2t)\|_{L^2(\mathbb R^3)}^2,
\]
and the right-hand side is finite for this fixed \(k\) by finite energy.
Thus each discarded index has absolute continuity of its \(L^3\)-integral at
the terminal slice.  Taking the maximum over this finite set gives the same
limit for the full supremum over \(k\).  This proves
\eqref{eq:auto-no-escape}.
\end{proof}

\begin{corollary}[Automatic admissibility of Type I concentration sequences]
\label{cor:auto-paperI-admissibility}
Let \((u,p,z_*)\), \(z_*=(x_*,T)\), be a finite-energy Type I terminal
singularity.  Then there exist \(\eta_v>0\) and \(\lambda_k\downarrow0\) such
that
\[
  C(z_*,\lambda_k)\ge\eta_v
  \qquad\text{for every }k,
\]
and the corresponding rescaled velocities satisfy the no-escape condition
\eqref{eq:auto-no-escape}.  Hence the concentration and no-escape
admissibility condition used in the companion Type I blow-up-limit paper is
automatic for finite-energy Type I singularities.
\end{corollary}

\begin{proof}
Since \(x_*\in\Sigma(T)\), \cref{cor:velocity-concentration} gives
\(\eta_v>0\) and \(\lambda_k\downarrow0\) with
\(C(z_*,\lambda_k)\ge\eta_v\) for all \(k\).  The no-escape condition for this
same sequence follows from \cref{thm:auto-velocity-no-escape}.
\end{proof}

\section{The Type I/Type II Blow-Up Rate Dichotomy}

Let \(x_0\in\Sigma(T)\), \(\eta>0\), and \(r_n\downarrow0\) be given by
\cref{thm:singular-positive}.  Define
\[
  u_n(y,s)=r_nu(x_0+r_ny,T+r_n^2s),
  \qquad
  p_n(y,s)=r_n^2p(x_0+r_ny,T+r_n^2s).
\]
Then \((u_n,p_n)\) are suitable weak solutions on expanding backward
cylinders, and by \cref{prop:critical-scaling} they retain a positive lower
bound for the mean-subtracted local scale-invariant quantity on a fixed
bounded cylinder.

\begin{definition}[Type I concentration sequence]
\label{def:typeI-concentration}
Let \((u,p)\) be a suitable weak solution on
\(\mathbb R^3\times(T-\delta,T)\), and let \(r_n\downarrow0\) be a positive
local concentration sequence at \((x_0,T)\).  The sequence is called Type I
if the concentration point \((x_0,T)\) satisfies a local Type I
scale-invariant bound:
there are \(\rho>0\) and \(M<\infty\), with \(\rho^2<\delta\), such that
\[
  \operatorname*{ess\,sup}_{T-\rho^2<t<T}
  \sqrt{T-t}\,\|u(t)\|_{L^\infty(B_\rho(x_0))}\le M.
\]
This is the Type I alternative.
\end{definition}

\begin{definition}[Local Type II alternative]
\label{def:typeII-alternative}
Let \((u,p)\) be a suitable weak solution on
\(\mathbb R^3\times(T-\delta,T)\), and let \(r_n\downarrow0\) be a positive
local concentration sequence at \((x_0,T)\).  The sequence belongs to the local
Type II alternative if the concentration point \((x_0,T)\) satisfies no
local Type I scale-invariant bound.  Equivalently, for every \(\rho>0\) with
\(\rho^2<\delta\),
\[
  \operatorname*{ess\,sup}_{T-\rho^2<t<T}
  \sqrt{T-t}\,\|u(t)\|_{L^\infty(B_\rho(x_0))}=\infty.
\]
\end{definition}

\begin{theorem}[Local blow-up rate dichotomy]\label{thm:typeI-typeII-dichotomy}
Let \((u,p)\) be a suitable weak solution on
\(\mathbb R^3\times(T-\delta,T)\), and suppose \(\Sigma(T)\ne\emptyset\).
Then there exist \(x_0\in\Sigma(T)\),
\(\eta>0\), and \(r_n\downarrow0\) such that
\[
  C((x_0,T),r_n)+D((x_0,T),r_n)\ge\eta
  \qquad\text{for all }n.
\]
The sequence \(r_n\) is a positive local concentration sequence.  The
corresponding concentration regime satisfies either the local Type I bound or
the local Type II alternative.  Consequently, once the vanishing case has been
excluded by \(\varepsilon\)-regularity, every finite-time singular case falls
under exactly one of the two blow-up rate alternatives.
\end{theorem}

\begin{proof}
If the local scale-invariant quantities vanished at every point of
\(\Sigma(T)\), then \cref{thm:no-concentration} would give
\(\Sigma(T)=\emptyset\), contrary to the hypothesis.  Therefore
\cref{thm:singular-positive} gives
\(x_0\in\Sigma(T)\), \(\eta>0\), and \(r_n\downarrow0\) with the stated lower
bound.  The rescaled sequence has a positive local scale-invariant lower bound
by \cref{prop:critical-scaling}.  By definition, either the Type I bound in
\cref{def:typeI-concentration} holds, or the local Type II alternative in
\cref{def:typeII-alternative} holds.
\end{proof}

\begin{remark}
This theorem is local.  Its conclusion is the existence of a positive
mean-subtracted scale-invariant concentration sequence and the dichotomy
between a local Type I bound and failure of every such bound at the
concentration point.
\end{remark}

\section{Local Pointwise Type I Bridge to the Terminal Concentration Package}
\label{sec:local-typeI-terminal-bridge-paper0}

We record the interface used by the conditional Type I proof setup
\cite{ProofSetupTypeI}.  The local pointwise Type I alternative in
\cref{def:typeI-concentration} is strong enough, after rescaling at the singular
point, to give compactness on every compact negative-time cylinder.  The
proof setup also records the endpoint weak-Serrin bridge:
for suitable weak solutions the local pointwise Type I envelope belongs to
\(L^{2,\infty}_tL^\infty_x\), hence the Albritton--Barker weak Serrin local
Type I package gives the scale-invariant \(A\)-bound on smaller cylinders and
therefore terminal velocity no-escape.  Thus positive local Type I
concentration sequences are admissible for the conditional Seregin extraction.

Throughout this section \(z_*=(x_*,T)\), and
\[
  Q_1(0,0)=B_1(0)\times(-1,0).
\]

\begin{definition}[Admissible local Type I bridge sequence]
\label{def:paper0-admissible-local-bridge-sequence}
For a local pointwise Type I point \(z_*=(x_*,T)\), a sequence
\(r_n\downarrow0\) is an admissible local Type I bridge sequence if the rescaled
velocities
\[
  u_n(x,t)=r_nu(x_*+r_nx,T+r_n^2t)
\]
satisfy, for some \(\eta_v>0\),
\[
   \int_{Q_1(0,0)} |u_n|^3\,dx\,dt\ge\eta_v
   \qquad\text{for every }n,
\]
and
\[
   \lim_{\sigma\downarrow0}
   \sup_n
   \int_{-\sigma}^{0}\int_{B_1}|u_n|^3\,dx\,dt=0 .
\]
\end{definition}

\begin{proposition}[Local pointwise Type I rescalings have terminal compactness]
\label{prop:paper0-local-typeI-terminal-compactness}
Let \((u,p)\) be a suitable weak solution on
\(\mathbb R^3\times(0,T)\) satisfying the finite-energy bounds
\[
  u\in L^\infty(0,T;L^2(\mathbb R^3))
  \cap L^2(0,T;\dot H^1(\mathbb R^3)).
\]
Let \(z_*=(x_*,T)\) be a singular point, and assume that \(z_*\) satisfies the
local pointwise Type I bound: there exist \(\rho>0\) and \(M<\infty\) such that
\begin{equation}\label{eq:paper0-local-typeI-bridge-bound}
  \operatorname*{ess\,sup}_{T-\rho^2<t<T}
  \sqrt{T-t}\,\|u(t)\|_{L^\infty(B_\rho(x_*))}\le M .
\end{equation}
For any sequence \(r_n\downarrow0\), define
\begin{equation}\label{eq:paper0-local-typeI-bridge-rescaling}
  u_n(x,t)=r_nu(x_*+r_nx,T+r_n^2t),
  \qquad
  p_n(x,t)=r_n^2p(x_*+r_nx,T+r_n^2t).
\end{equation}
Then the following hold.
\begin{enumerate}[label=\textup{(\roman*)}]
\item If
\[
  K=B_R\times[-S,-\sigma]\Subset\mathbb R^3\times(-\infty,0),
  \qquad R,S<\infty,\quad \sigma>0,
\]
then, for all sufficiently large \(n\),
\begin{equation}\label{eq:paper0-local-typeI-bridge-Linfty}
  \|u_n\|_{L^\infty(K)}\le M\sigma^{-1/2}.
\end{equation}
\item On each such compact cylinder \(K\), after subtracting time-dependent
local pressure gauges \(a_{n,K}(t)\), the sequence satisfies uniform local
Seregin bounds
\begin{equation}\label{eq:paper0-local-typeI-bridge-Seregin-bounds}
  \|u_n\|_{L^\infty_tL^2_x(K)}
  +\|\nabla u_n\|_{L^2(K)}
  +\|p_n-a_{n,K}(t)\|_{L^{3/2}(K)}
  \le C_K .
\end{equation}
Here \(C_K\) may depend on \(K,M,\rho\), and on the finite Leray energy level
of the original solution, but it does not depend on \(n\) and does not use a
terminal no-escape assumption.
\item After passing to a subsequence, \((u_n,p_n-a_{n,K})\) converges on every
compact \(K\Subset\mathbb R^3\times(-\infty,0)\) to an ancient suitable weak
solution \((U,Q)\), with
\[
  u_n\to U\quad\text{strongly in }L^q(K),\qquad 1\le q<\infty,
\]
\[
  p_n-a_{n,K}(t)\rightharpoonup Q
  \quad\text{weakly in }L^{3/2}(K),
\]
and
\[
  |U(x,t)|\le M(-t)^{-1/2},\qquad t<0.
\]
\end{enumerate}
\end{proposition}

\begin{proof}
Fix \(K=B_R\times[-S,-\sigma]\Subset\mathbb R^3\times(-\infty,0)\).  For
large \(n\),
\[
  r_nR<\rho/2,\qquad r_n^2S<\rho^2/2.
\]
Thus \(x_*+r_nx\in B_\rho(x_*)\) and
\(T+r_n^2t\in(T-\rho^2,T)\) for a.e. \((x,t)\in K\).  The local pointwise Type I bound
\eqref{eq:paper0-local-typeI-bridge-bound} gives
\[
  |u_n(x,t)|
  =r_n|u(x_*+r_nx,T+r_n^2t)|
  \le r_nM(T-(T+r_n^2t))^{-1/2}
  =M(-t)^{-1/2}
  \le M\sigma^{-1/2}.
\]
This proves \eqref{eq:paper0-local-typeI-bridge-Linfty}.  It also gives
uniform local \(L^\infty_tL^2_x\) and \(L^3_{x,t}\) bounds on \(K\).

It remains to explain the pressure and dissipation bounds.  Use the
whole-space Leray pressure representative, which is legitimate in the
finite-energy class up to addition of functions of time.  Choose
\(\chi_R\in C_c^\infty(B_{8R})\), \(\chi_R\equiv1\) on \(B_{4R}\), and write
inside \(B_{4R}\)
\[
  p_n=q_n+h_n,\qquad
  q_n=\mathcal R_i\mathcal R_j(\chi_Ru_{n,i}u_{n,j}).
\]
Calderon--Zygmund boundedness and the local \(L^3\)-bound for \(u_n\) give a
uniform \(L^{3/2}\)-bound for \(q_n\).  For \(h_n\), take the gauge
\(a_{n,K}(t)=h_n(0,t)\).  The harmonic pressure kernel estimate gives
\[
  |h_n(x,t)-a_{n,K}(t)|
  \le
  C_R\sum_{m\ge2}2^{-4m}
  \int_{2^mR<|z|<2^{m+1}R}|u_n(z,t)|^2\,dz,
  \qquad x\in B_{2R}.
\]
For annuli whose physical image remains inside \(B_\rho(x_*)\), the local
Type I bound controls the \(L^2\)-mass by
\[
  C M^2\sigma^{-1}(2^mR)^3,
\]
and the resulting series is geometric.  For the outer annuli, the finite
Leray energy gives
\[
  \int_{\mathbb R^3}|u_n(z,t)|^2\,dz
  =r_n^{-1}\int_{\mathbb R^3}|u(y,T+r_n^2t)|^2\,dy,
\]
while the first outer annulus has radius comparable to \(\rho/r_n\).  Hence
the tail of the kernel series is bounded by
\[
  C_{\rho,R}\,r_n^3\|u\|_{L^\infty(0,T;L^2)}^2,
\]
which is uniform in \(n\).  This proves the pressure-gauge part of
\eqref{eq:paper0-local-typeI-bridge-Seregin-bounds}.

Apply the local energy inequality for \((u_n,p_n)\) with a cutoff equal to one
on \(K\) and supported in a slightly larger compact cylinder.  The local
kinetic, cubic velocity, and pressure-flux terms are bounded by the estimates
just proved, so \(\|\nabla u_n\|_{L^2(K)}\le C_K\).  The Navier--Stokes
equation then bounds \(\partial_tu_n\) locally in a negative Sobolev space.
Aubin--Lions compactness and the pressure bound give the asserted Seregin
subsequence convergence.  Passing to the limit in the equations and local
energy inequality gives an ancient suitable weak solution, and the pointwise
Type I bound passes to the limit.
\end{proof}

\begin{corollary}[Local pointwise Type I bridge sequences are admissible]
\label{cor:paper0-local-typeI-bridge-admissible}
Under the hypotheses of
\cref{prop:paper0-local-typeI-terminal-compactness}, let \(r_n\downarrow0\) be
any sequence such that the rescaled velocities satisfy
\[
  \int_{Q_1(0,0)} |u_n|^3\,dx\,dt\ge\eta_v>0
  \qquad\text{for all }n .
\]
Then \(r_n\) is an admissible local Type I bridge sequence in the sense of
\cref{def:paper0-admissible-local-bridge-sequence}.
\end{corollary}

\begin{proof}
The displayed positive concentration lower bound is the first requirement in
\cref{def:paper0-admissible-local-bridge-sequence}.  The terminal
no-escape statement is taken here as the endpoint weak-Serrin bridge from the
conditional proof setup \cite{ProofSetupTypeI}; analytically, it is the
application of the Albritton--Barker weak Serrin local Type I estimate
\cite{AlbrittonBarker2019} to the local pointwise Type I envelope.  This is
exactly the second requirement in
\cref{def:paper0-admissible-local-bridge-sequence}.
\end{proof}

\begin{lemma}[Local Type I concentration enters the terminal dichotomy]
\label{lem:paper0-local-typeI-into-terminal-dichotomy}
Under the hypotheses of
\cref{prop:paper0-local-typeI-terminal-compactness}, let \(r_n\downarrow0\) be
a velocity concentration sequence satisfying
\[
  \int_{Q_1(0,0)} |u_n|^3\,dx\,dt\ge\eta_v>0
  \qquad\text{for all }n .
\]
Then \(r_n\) is admissible by
\cref{cor:paper0-local-typeI-bridge-admissible}.  There are
\(\sigma_0\in(0,1)\) and \(\eta_1>0\) such that
\begin{equation}\label{eq:paper0-local-typeI-bridge-away-mass}
  \int_{-1}^{-\sigma_0}\int_{B_1}|u_n|^3\,dx\,dt\ge\eta_1
  \qquad\text{for all }n.
\end{equation}
Consequently, after applying the terminal concentration package stated in the
conditional proof setup \cite{ProofSetupTypeI} with a threshold
\(\varepsilon_{\rm sd}\le\eta_1\), the retained compact concentration issued
from the origin recentering has no loss of local compactness.  The paired
terminal decomposition produces an active retained profile locus plus a
residual remainder.  With realized critical-tail stratification and active
path-space recurrence assumed as part of the proof setup, the terminal package
excludes mixed compact--diffuse remainders,
compact--noncompact remainders, and infinite active descendant branches.  Thus
the retained compact concentration is reduced to a single retained
concentrating profile or a finite separated family of retained concentrating
profiles.
\end{lemma}

\begin{proof}
The no-escape condition in
\cref{def:paper0-admissible-local-bridge-sequence} gives \(\sigma_0>0\) such
that the terminal layer \(B_1\times(-\sigma_0,0)\) carries at most half of the
fixed \(Q_1\) velocity mass.  This proves
\eqref{eq:paper0-local-typeI-bridge-away-mass}.  The compactness part follows
from \cref{prop:paper0-local-typeI-terminal-compactness} on the compact
negative-time cylinder \(B_1\times[-1,-\sigma_0]\).  Strong local \(L^3\)
convergence on this cylinder passes the positive mass to the ancient limit.
The local CKN measure in the terminal package dominates the velocity
contribution, so after a finite covering of
\(B_1\times(-1,-\sigma_0)\) by unit terminal cylinders, one origin-frame
recentering is \(\varepsilon_{\rm sd}\)-concentrating for a fixed
\(\varepsilon_{\rm sd}>0\) depending on \(\eta_1\).  This is the only point at
which terminal no-escape is used.

The terminal exhaustion statement used here is the following conditional input
from the proof setup: the active retained profile locus and residual remainder
are paired; compact core plus diffuse or noncompact remainders are routed to
the realized critical-tail rigidity theorem; and infinite active cascades are
ruled out by recurrence.  Therefore only a single retained profile or a finite
separated retained family remains at this bridge stage.
\end{proof}

\begin{theorem}[Local Type I entry into the full exclusion pipeline]
\label{thm:paper0-local-typeI-entry-pipeline}
Let \((u,p)\) be a finite-energy suitable weak solution
on \(\mathbb R^3\times(0,T)\), and let \(z_*=(x_*,T)\) be a singular point
satisfying the local pointwise Type I bound
\eqref{eq:paper0-local-typeI-bridge-bound}.  Then the rescalings at \(z_*\)
generate a bounded centered Type I Seregin profile with retained local
concentration, and the conditional terminal package applies to that profile.  In
particular, no finite-energy local pointwise Type I singular point exists in
the completed Papers~0--IV pipeline.
\end{theorem}

\begin{proof}
Choose the positive \(Q_1\) velocity concentration sequence supplied by the
local concentration theorem.  By
\cref{cor:paper0-local-typeI-bridge-admissible}, this sequence satisfies
terminal no-escape and is admissible.
By \cref{lem:paper0-local-typeI-into-terminal-dichotomy}, the local Type I
singularity generates retained compact concentration with local compactness.
The compactness statement gives an ancient Type I limit satisfying
\(|U(x,t)|\le M(-t)^{-1/2}\); after the centered change of variables this is a
bounded centered Type I Seregin profile.  If all retained profiles generated by
the terminal package belong to lower classes already closed in Papers I--III,
then the contradiction is supplied by those earlier class exclusions.  If some
retained profile lies outside the lower classes, it is a residual candidate.
For residual candidates, we assume the terminal local stratification package
from the proof setup \cite{ProofSetupTypeI}: finite separated retained
families are excluded, infinite active branches are excluded by active
path-space recurrence, mixed compact--diffuse tails are routed through
critical-tail compactification and realized critical-tail rigidity, and a
remaining single retained profile is forced through indecomposability,
parasitic-free mildness inheritance, and the endpoint Liouville contradiction.
Thus the assumed local Type I singular point cannot exist in the completed
pipeline.
\end{proof}

\begin{corollary}[Local Type I alternatives are excluded]
\label{cor:paper0-local-typeII-after-bridge}
Under the conditional proof-setup hypotheses, the local Type I alternative is
excluded for every finite-energy singular point.  Hence every finite-energy
singular point belongs to the local Type II alternative in the dichotomy above.
\end{corollary}

\begin{proof}
This is \cref{thm:paper0-local-typeI-entry-pipeline} restated in the language
of the local Type I/local Type II dichotomy.
\end{proof}

\begin{thebibliography}{9}

\bibitem{AlbrittonBarker2019}
D.~Albritton and T.~Barker,
\emph{On local Type I singularities of the Navier--Stokes equations and
Liouville theorems},
J. Math. Fluid Mech. \textbf{21} (2019), Article 43.

\bibitem{CKN1982}
L.~Caffarelli, R.~Kohn, and L.~Nirenberg,
\emph{Partial regularity of suitable weak solutions of the Navier--Stokes
 equations},
Comm. Pure Appl. Math. \textbf{35} (1982), no. 6, 771--831,
doi:\ \nolinkurl{10.1002/cpa.3160350604}.

\bibitem{ESS2003}
L.~Escauriaza, G.~A.~Seregin, and V.~{\v S}verak,
\emph{\(L_{3,\infty}\)-solutions of the Navier--Stokes equations and
 backward uniqueness},
Russian Math. Surveys \textbf{58} (2003), no. 2, 211--250,
doi:\ \nolinkurl{10.1070/RM2003v058n02ABEH000609}.

\bibitem{ProofSetupTypeI}
G.~Duran Ballester,
\emph{Conditional Type I Navier--Stokes proof setup},
\texttt{docs/source/type\_1/unconditional\_typeI\_chain/proof\_setup.tex}.

\bibitem{Giga1986}
Y.~Giga,
\emph{Solutions for semilinear parabolic equations in \(L^p\) and regularity
 of weak solutions of the Navier--Stokes system},
J. Differential Equations \textbf{62} (1986), no. 2, 186--212,
doi:\ \nolinkurl{10.1016/0022-0396(86)90096-3}.

\bibitem{Leray1934}
J.~Leray,
\emph{Sur le mouvement d'un liquide visqueux emplissant l'espace},
Acta Math. \textbf{63} (1934), no. 1, 193--248,
doi:\ \nolinkurl{10.1007/BF02547354}.

\bibitem{Lin1998}
F.-H.~Lin,
\emph{A new proof of the Caffarelli--Kohn--Nirenberg theorem},
Comm. Pure Appl. Math. \textbf{51} (1998), no. 3, 241--257,
doi:\ \nolinkurl{10.1002/(SICI)1097-0312(199803)51:3<241::AID-CPA2>3.0.CO;2-A}.

\bibitem{Scheffer1976}
V.~Scheffer,
\emph{Partial regularity of solutions to the Navier--Stokes equations},
Pacific J. Math. \textbf{66} (1976), no. 2, 535--552,
doi:\ \nolinkurl{10.2140/pjm.1976.66.535}.

\bibitem{Seregin2012}
G.~Seregin,
\emph{A certain necessary condition of potential blow up for Navier--Stokes
 equations},
Comm. Math. Phys. \textbf{312} (2012), no. 3, 833--845,
doi:\ \nolinkurl{10.1007/s00220-011-1391-x}.

\bibitem{Serrin1962}
J.~Serrin,
\emph{On the interior regularity of weak solutions of the Navier--Stokes
 equations},
Arch. Rational Mech. Anal. \textbf{9} (1962), 187--195,
doi:\ \nolinkurl{10.1007/BF00253344}.

\bibitem{Stein1970}
E.~M. Stein,
\emph{Singular Integrals and Differentiability Properties of Functions},
Princeton Mathematical Series, No.~30, Princeton University Press, 1970.

\end{thebibliography}

\end{document}
